\section{Discussion and Limitations}\label{sec:disc}

\noindent\textbf{What the evidence shows.} Across E1--E4 the two backends behave as one engine wearing two geometries. E1's near-identical F1 trajectories and E4's live-cell counts are the empirical face of \eqref{eq:objective}: \code{LayeredMap} and \code{PeriodicityModel}, exercised directly in E2 and E3, never see which backend produced the id they update. The two places where the backends do diverge---the high-clutter precision gap and the per-call latency gap---trace to native sampling and to the 3D ray march, respectively, rather than to the shared classifier.

\noindent\textbf{Cost Analysis.} The window update is linear in the number of live cells (\S\ref{subsec:classify}), so cost is governed by how many cells a backend creates. The flat voxel hash is the source of \code{voxel3d}'s 4.4--10.0$\times$ larger live-cell count: every sampled free voxel becomes its own entry at 56~B, with no spatial compression. OctoMap \citep{hornung2013octomap} and UFOMap \citep{duberg2020ufomap} collapse coherent free or occupied regions into single octree nodes and would scale better to very large, sparsely occupied 3D extents. We accept the trade for $O(1)$ insert and lookup and a backend simple enough to unit-test without ROS; even the largest configuration tested needs 37.0~MB.

\noindent\textbf{Forgetting outside the field of view.} Remark~\ref{rem:fov} is an operational limitation, not a footnote. Decay runs only inside the touched branch, so a transient object that accumulates strong evidence and then leaves sensor coverage for good---a box dragged into a corridor the robot no longer visits---keeps its evidence indefinitely. Forgetting in \sname\ is coupled to re-observation, not to elapsed time. In a bounded, frequently re-swept environment, the setting E1--E4 implicitly assume, this rarely matters; in a sparsely revisited one it is a real failure mode that the current code does not address.

\noindent\textbf{Design trade-offs.} Three simplifications are deliberate. A single global \code{survival\_decay} replaces the per-cell survival posterior of \citet{rosen2016persistence} and the learned transition rates of \citet{tipaldi2013lifelong} and \citet{saarinen2012imac}; E3 shows that, once paired with hysteresis, one uniform rate already yields F1 1.000 at both $\lambda{=}0.90$ and $1.0$, but a scene with sharply heterogeneous per-cell dynamics is untested. The flat hash trades memory scaling for simplicity, as above. And two backends behind an \code{if}/\code{else} factory is the right amount of mechanism for two, but a third backend would argue for a registry such as \code{pluginlib} \citep{ros2pluginlib}; this is a design judgement, not a measured result.

\noindent\textbf{When not to use \sname.} It is not a substitute for SLAM: it estimates no pose and closes no loops, so a deployment without an external localizer publishing the \code{map}$\to$sensor transform has no path to a map. It is not multi-session mapping: comparing Monday's map with Friday's or replaying a past session needs the versioning of LT-mapper or \citet{yang2025lifelong3d}. And it has no object-level reasoning, unlike Khronos's object-level scene graph \citep{schmid2024khronos}. Within those bounds---one continuous session behind an external localizer, at the raw cell level, in 2D or 3D---it is the intended fit.

\noindent\textbf{Limitations.} All evidence is synthetic: every number comes from a deterministic seeded harness driving hand-authored hit and miss patterns through the shipped engine, with no real sensor, robot motion or field deployment, whereas ELite, LT-mapper and \citet{berrio2021longterm} report multi-session field evidence. E1--E4 are single-seed and carry no dispersion estimate; only E5 is multi-seed. No external baseline is run. E4's microsecond figures come from one unrecorded machine. The engine models single-session persistence only, forgets only what it re-observes (Remark~\ref{rem:fov}), and loses low-duty-cycle periodic cells to the prune/maturity race of Proposition~\ref{prop:race}. The periodicity test has limits of its own. What is proved (Propositions~\ref{prop:chernoff} and~\ref{prop:spend}) is a per-history bound under an iid null and occupancy-independent touches: at the shipped level, a cell that is never pruned is labelled Periodic at some window of the run with probability at most 0.2. The 0.01 target is a measurement: at most 2 of 2000 clutter cells within 1024 windows on held-out seeds (E5), and for pruned clutter the proved bound grows with the number of re-created histories. Temporally correlated clutter violates the null, and at the shipped $T{=}24$ the test labels most slowly switching cells Periodic (1908 of 2000 in E5). Spending costs detection delay: a noiseless 50\%-duty door with $T{=}8$ needs 25 windows instead of 15, and one with 10\% flipped windows 48 instead of 26. Finally, the four-class confusion shows that the dominant labelling error under the E2 parameters is not periodic at all: 1221 of 2000 aperiodic cells end falsely Static, which is a property of the graduation band without decay. The natural next steps follow directly: a null that allows temporal correlation, for example by testing the harmonic against a local noise floor rather than against white noise, and an anytime-valid test that also covers re-created histories; field validation on a 2D wheeled platform and a 3D-LiDAR platform; per-cell learned decay; exempting recently active cells from pruning, or feeding a running amplitude estimate into \eqref{eq:prune}; and relaxing untouched cells toward the prior on a tick basis, without changing the touched-cell path that E1--E4 characterise.
